% Part: incompleteness
% Chapter: incompleteness-provability
% Section: introduction

\documentclass[../../../include/open-logic-section]{subfiles}

\begin{document}

\olfileid[hi]{inc}{inp}{int}

\olsection{परिचय}

हिल्बर्ट का विचार था कि प्राकृतिक संख्याओं जैसी किसी गणितीय संरचना के लिए
अभिगृहीतों का तंत्र तब तक अपर्याप्त है, जब तक वह उस संरचना के बारे में सभी
सत्य कथनों को व्युत्पन्न न करने दे। निगमन के औपचारिक तंत्रों में उनकी बाद
की रुचि के साथ मिलाकर इससे संकेत मिलता है कि उनके विचार में हमें यह
सुनिश्चित करना चाहिए कि, मान लें, प्राकृतिक संख्याओं के बारे में तर्क करने
के लिए प्रयुक्त औपचारिक तंत्र न केवल संगत हो, बल्कि \emph{पूर्ण} भी हो;
अर्थात् उसकी भाषा का प्रत्येक कथन या तो !!{derivable} हो या उसका निषेध
हो। ग्योडेल (G\"odel) का प्रथम अपूर्णता-प्रमेय दिखाता है कि
अभिगृहीतों का ऐसा कोई तंत्र नहीं है: अंकगणित का कोई पूर्ण, संगत,
!!{axiomatizable} औपचारिक तंत्र नहीं है। वास्तव में कोई भी ``पर्याप्त
प्रबल,'' संगत, !!{axiomatizable} गणितीय सिद्धांत पूर्ण नहीं है।

हिल्बर्ट का अधिक महत्त्वपूर्ण लक्ष्य---आधुनिक (``चिरसम्मत'') गणित के
औचित्य के उनके कार्यक्रम का केंद्रीय अंश---चिरसम्मत तर्क का निरूपण करने
वाले औपचारिक तंत्रों के लिए परिमितीय संगतता-प्रमाण खोजना था। इसलिए हिल्बर्ट
के कार्यक्रम के संदर्भ में ग्योडेल (G\"odel) का द्वितीय अपूर्णता-प्रमेय कहीं अधिक
बड़ा आघात था। द्वितीय अपूर्णता-प्रमेय को प्रथम अपूर्णता-प्रमेय की तरह कुछ
अस्पष्ट शब्दों में कहा जा सकता है। मोटे तौर पर वह कहता है कि अंकगणित का
कोई पर्याप्त प्रबल सिद्धांत अपनी संगतता स्वयं सिद्ध नहीं कर सकता। यहाँ
``पर्याप्त प्रबल'' में हमें~$\Th{Q}$ से कुछ अधिक सम्मिलित करना होगा।

ग्योडेल (G\"odel) के अपूर्णता-प्रमेय के मूल प्रमाण के पीछे का विचार एपिमेनिडीज़
(Epimenides) विरोधाभास में मिलता है। क्रीट निवासी एपिमेनिडीज़ ने कहा था कि
सभी क्रीट निवासी झूठे हैं; इस विरोधाभास का अधिक सीधा रूप है कथन ``यह वाक्य
असत्य है।'' मूलतः सत्यता के स्थान पर !!{derivability} रखकर ग्योडेल (G\"odel) ऐसे
!!a{sentence} को औपचारिक बनाने में सफल हुए जो घुमावदार ढंग से कहता है कि
वह स्वयं !!{derivable} नहीं है। यदि वह !!{sentence},~!!{derivable} होता,
तो सिद्धांत असंगत होता। ग्योडेल (G\"odel) ने यह भी दिखाया कि जिस अभिगृहीत-तंत्र पर
वे विचार कर रहे थे, उससे उस !!{sentence} का निषेध भी !!{derivable} नहीं
है। (इस दूसरे भाग के लिए ग्योडेल (G\"odel) को यह मानना पड़ा कि सिद्धांत~$\Th{T}$
``$\omega$-संगत'' है। $\omega$-संगतता, संगतता से संबंधित किंतु उससे अधिक
प्रबल गुण है।\footnote{अर्थात् प्रत्येक $\omega$-संगत सिद्धांत संगत है,
किंतु विलोम सत्य नहीं है।} ग्योडेल (G\"odel) के कुछ वर्ष बाद रॉसर (Rosser) ने दिखाया
कि~$\Th{T}$ की केवल साधारण संगतता मानना पर्याप्त है।)

पहली चुनौती यह समझना है कि स्वयं को निर्दिष्ट करने वाला कोई !!a{sentence}
कैसे बनाया जा सकता है। $\Th{Q}$ की भाषा के प्रत्येक !!{formula}~$!A$ के
लिए~$\gn{!A}$ को~$\Gn{!A}$ के संगत अंक-सूचक मानें। इसका अर्थ ध्यान से
सोचें: $!A$,~$\Th{Q}$ की भाषा का !!a{formula} है; $\Gn{!A}$ कोई प्राकृतिक
संख्या है; और $\gn{!A}$,~$\Th{Q}$ की भाषा का कोई \emph{पद} है। अतः
$\Th{Q}$ की भाषा के प्रत्येक !!{formula}~$!A$ का एक \emph{नाम},~$\gn{!A}$,
है, जो~$\Th{Q}$ की भाषा का पद है; इससे हमें ऐसा संकल्पनात्मक ढाँचा मिलता
है जिसमें~$\Th{Q}$ की भाषा के !!{formula}s दूसरे !!{formula}s के बारे में
बातें ``कह'' सकते हैं। अगली उपपत्ति नियत-बिंदु उपपत्ति कहलाती है।

\begin{lem}
$\Th{T}$,~$\Th{Q}$ का कोई विस्तार हो, और $!B(x)$ ऐसा कोई !!{formula} हो
जिसमें केवल चर~$x$ मुक्त है। तब कोई !!a{sentence}~$!A$ ऐसा है कि
$\Th{T} \Proves!A \liff !B(\gn{!A})$।
\end{lem}

उपपत्ति कहती है कि कोई गुण $!B(x)$ दिए जाने पर ऐसा !!a{sentence}~$!A$ है
जो कहता है ``$!B(x)$ मेरे बारे में सत्य है,'' और $\Th{T}$ इसे ``जानता''
है।

ऐसा !!a{sentence} हम कैसे बना सकते हैं? क्वाइन (Quine) द्वारा दिए गए
एपिमेनिडीज़ विरोधाभास के निम्न रूप पर विचार करें:
\begin{quote}
``अपना उद्धरण पहले रखने पर असत्य देता है'' अपना उद्धरण पहले रखने पर असत्य
देता है।
\end{quote}
यह वाक्य सीधे स्वयं को निर्दिष्ट नहीं करता। वह केवल उद्धरण-चिह्नों के बीच
की वाक्यविन्यासिक वस्तुओं के बारे में दावा करता है और ऐसा करते हुए निम्न
वाक्यों के समान स्तर पर है:
\begin{enumerate}
\item ``रॉबर्ट'' अच्छा नाम है।
\item ``मैं दौड़ा।'' छोटा वाक्य है।
\item ``तीन शब्द हैं'' में तीन शब्द हैं।
\end{enumerate}
किंतु जब कोई ``अपना उद्धरण पहले रखने पर असत्य देता है'' पदबंध को लेकर उसके
पहले उसी का उद्धृत रूप रखता है, तो क्या होता है? तब मूल वाक्य ही मिल जाता
है!{} संक्षेप में, वाक्य कहता है कि वह असत्य है।

\end{document}
